% !TEX root = ../math.tex
\chapter*{Introduction}
\phantomsection
\addcontentsline{toc}{chapter}{Introduction}

% TODO: Add introduction text.

Dotter is a text-entry technology that allows users to type over a 1-bit channel.
It allows a simple input mechanism such as blinking or pressing a single button
to be used to type.
Dotter milks as much information out of that 1-bit channel as possible,
allowing an experienced user to type standard english text at speeds of 20 words per minute
only by blinking.

\begin{quote}
\textbf{Dotter only has one rule:} Identify the longest visible prefix of your target text and synchronize your signal to its timer.
\end{quote}

Dotter was inspired by the project ``Dasher''~\cite{Ward2000Dasher}, a text-entry technology
developed by David MacKay's group at the University of Cambridge that took inspiration from arithmetic coding to
map the space of all possible text sequences to the unit interval. Assuming a version
of Fitt's law, Dasher is an optimal text-entry technology for 1-dimensional continuous (pointing) input.

This document is a mathematical analysis of Dotter as well as an implementation guide,
describing the architecture and key algorithms of Dotter.
Chapter \ref{chap:architecture} describes the architecture of Dotter,
outlining its key data structures and algorithms.
The succeeding chapters describe individual components of Dotter.
The final chapter gives theory and advice on how to learn to type with Dotter.
(If the person reading this only wants to learn how to use Dotter, they should skip to the final chapter.)

The name ``Dotter'' is a nod the most famous 1-bit text-entry system, telegraphy,
system which consisted of ``dots'' and ``dashes.''
The name ``Dasher'' was already taken, so I had to go with this.
